\section{Thiết kế Knowledge Graph}

Knowledge Graph là lớp mở rộng giúp RAG không chỉ truy hồi đoạn văn gần nghĩa, mà còn hiểu quan hệ giữa các thực thể. Trong bài toán tuyển sinh, các entity quan trọng gồm ngành học, mã ngành, tổ hợp môn, phương thức xét tuyển, điểm chuẩn, học phí, năm tuyển sinh, cơ sở đào tạo và loại tài liệu.

Sau khi chunk được chuẩn hóa, hệ thống có thể trích xuất entity, chuẩn hóa entity, trích xuất relation, xử lý xung đột và ghi vào graph store. Các entity và relation này được liên kết ngược về chunk và tài liệu gốc để bảo đảm khả năng truy vết.

\begin{figure}[H]
    \centering
    \includegraphics[width=0.9\textwidth]{\detokenize{content/source/RAG/Pipeline xây dựng Knowledge Graph từ chunk đã chuẩn hóa.png}}
    \caption{Pipeline xây dựng Knowledge Graph từ chunk đã chuẩn hóa}
    \label{fig:rag-knowledge-graph-pipeline}
\end{figure}

Ví dụ entity:

\begin{itemize}
    \item Ngành: Công nghệ thông tin.
    \item Mã ngành: 7480201KMA.
    \item Tổ hợp: A00, A01, D01.
    \item Phương thức: xét điểm thi THPT, xét học bạ.
    \item Năm: 2025, 2026.
    \item Tài liệu: Quy chế tuyển sinh 2026.
\end{itemize}

Ví dụ relation:

\begin{itemize}
    \item Công nghệ thông tin -- có mã ngành -- 7480201KMA.
    \item Công nghệ thông tin -- xét tổ hợp -- A00.
    \item Công nghệ thông tin -- xét tổ hợp -- A01.
    \item Công nghệ thông tin -- xét tổ hợp -- D01.
    \item Công nghệ thông tin -- có điểm chuẩn năm 2025 -- 24.42.
    \item Quy chế tuyển sinh 2026 -- quy định -- phương thức xét tuyển.
\end{itemize}

Flow report mô tả rõ giá trị của Knowledge Graph: vector search trả lời tốt câu hỏi gần nghĩa, nhưng graph giúp trả lời các câu hỏi có quan hệ, có điều kiện và cần tổng hợp; ví dụ thay vì chỉ tìm đoạn nói về ``D01'', hệ thống có thể tìm các ngành có relation với tổ hợp D01, sau đó lấy chunk nguồn để trích dẫn.

Thiết kế schema rút gọn:

\begin{table}[H]
\centering
\begin{tabular}{|p{3.8cm}|p{9cm}|}
\hline
\textbf{Đối tượng} & \textbf{Thuộc tính} \\
\hline
Entity & \texttt{entity\_id}, \texttt{canonical\_name}, \texttt{aliases}, \texttt{entity\_type}, \texttt{description}, \texttt{source\_chunk\_id}, \texttt{document\_id}, \texttt{status}. \\
\hline
Relation & \texttt{relation\_id}, \texttt{source\_entity}, \texttt{target\_entity}, \texttt{relation\_type}, \texttt{description}, \texttt{source\_chunk\_id}, \texttt{confidence}, \texttt{status}. \\
\hline
\end{tabular}
\caption{Schema rút gọn cho entity và relation trong Knowledge Graph}
\label{tab:rag-kg-schema}
\end{table}

Knowledge Graph đặc biệt hữu ích với các câu hỏi như: những ngành nào xét D01, ngành Công nghệ thông tin có xét học bạ không, mã ngành 7480201 thuộc ngành nào, ngành An toàn thông tin có những tổ hợp nào hoặc điểm chuẩn ngành này qua các năm thay đổi thế nào.
